\section{Method}
\label{sec:method}

\subsection{Problem Formulation}

Let $(\bm{x}_i,y_i)$ be an image and its label. A frozen teacher $f_T$ and a
trainable student $f_S$ produce logits $\bm{z}_i^T$ and $\bm{z}_i^S$. The student
is trained with the total objective

\begin{equation}
\mathcal{L}
=
\mathcal{L}_{\mathrm{CE}}
+ \lambda_{\mathrm{KD}}\mathcal{L}_{\mathrm{KD}}
+ \lambda_{\mathrm{geo}}\mathcal{L}_{\mathrm{geo}}.
\label{eq:total-loss}
\end{equation}

The first term is ordinary cross-entropy:

\begin{equation}
\mathcal{L}_{\mathrm{CE}}
= -\frac{1}{N}\sum_{i=1}^{N}\log p_S(y_i\mid\bm{x}_i).
\end{equation}

For output-level transfer, softened probabilities at temperature $\tau$ are

\begin{equation}
q_{i,c}^{M}(\tau)
= \frac{\exp(z_{i,c}^{M}/\tau)}
{\sum_{k}\exp(z_{i,k}^{M}/\tau)}, \qquad M\in\{T,S\},
\end{equation}

and the distillation loss is

\begin{equation}
\mathcal{L}_{\mathrm{KD}}
= \frac{\tau^2}{N}\sum_{i=1}^{N}
\operatorname{KL}\!\left(
\bm{q}_{i}^{T}(\tau)\,\|\,\bm{q}_{i}^{S}(\tau)
\right).
\end{equation}

\subsection{Spatial Geometry Transfer}

At a selected layer, let $\bm{a}_{i,p}^{M}\in\mathbb{R}^{d}$ denote the aligned
feature vector at spatial position $p$. We first normalize it:

\begin{equation}
\widehat{\bm{a}}_{i,p}^{M}
= \frac{\bm{a}_{i,p}^{M}}
{\lVert\bm{a}_{i,p}^{M}\rVert_2+\varepsilon}.
\end{equation}

The cosine relation between positions $p$ and $q$ is
$r_{i,pq}^{M}=(\widehat{\bm{a}}_{i,p}^{M})^{\top}
\widehat{\bm{a}}_{i,q}^{M}$. The geometry loss matches these relations:

\begin{equation}
\mathcal{L}_{\mathrm{geo}}
= \frac{1}{N|\mathcal{P}|}
\sum_{i=1}^{N}\sum_{(p,q)\in\mathcal{P}}
\left(r_{i,pq}^{S}-r_{i,pq}^{T}\right)^2.
\label{eq:geometry-loss}
\end{equation}

Unlike direct feature regression, \cref{eq:geometry-loss} preserves how spatial
positions relate to one another while allowing teacher and student features to
use different coordinate systems. The complete objective is given in
\cref{eq:total-loss}.
